\documentclass{article}
\usepackage[T1]{fontenc}
\usepackage{lmodern}
\usepackage{graphicx}
\usepackage{amsmath,amssymb}
\usepackage[a4paper,margin=2cm]{geometry}
\usepackage[absolute,overlay]{textpos}
\setlength{\TPHorizModule}{1mm}
\setlength{\TPVertModule}{1mm}
\usepackage[indonesian]{babel}
\usepackage{hyperref}
\setlength{\emergencystretch}{2em}
% unit: Halaman 14

\begin{document}

% === Halaman 14 (python/native) ===

\#define LENGTH 16            /* Jumlah byte di dalam blok atau kunci */ 

\#define NROWS 4              /* Jumlah baris di dalam state */ 

\#define NCOLS 4              /* Jumlah kolom di dalam state */ 

\#define ROUNDS 10            /* Jumlah putaran */ 

typedef unsigned char byte;  /* unsigned 8-bit integer */ 


rijndael (byte plaintext[LENGTH], byte ciphertext[LENGTH],  

          byte key[LENGTH]) 

\{ 

 int r;                     /* pencacah pengulangan */ 

 byte state[NROWS][NCOLS];  /* state sekarang */ 

 struct\{byte k[NROWS][NCOLS];\} rk[ROUNDS + 1];  /* kunci pada setiap putaran */ 


 KeyExpansion(key, rk);       /* bangkitkan kunci setiap putaran */ 

 CopyPlaintextToState(state, plaintext);  /* inisialisasi 

                                             state sekarang */ 

 AddRoundKey(state, rk[0]);  /* XOR key ke dalam state */ 


 for (r = 1; r<= ROUNDS - 1; r++) 

 \{ 

  SubBytes(state);        /* substitusi setiap byte dengan S-box */  

  ShiftRows(state);       /* rotasikan baris i sejauh i byte */   

  MixColumns(state);      /* acak masing-masing kolom */ 

  AddRoundKey(state, rk[r]);   /* XOR key ke dalam state */ 

 \} 

 SubBytes(state);        /* substitusi setiap byte dengan S-box */  

 ShiftRows(state);       /* rotasikan baris i sejauh i byte */   

 AddRoundKey(state, rk[ROUNDS]);   /* XOR key ke dalam state */ 


 CopyStateToCiphertext(ciphertext, state);  /* blok cipherteks yang dihasilkan */  


\}

\end{document}
